\documentclass[tikz,border=8pt]{standalone}
\usepackage{fontspec}
\usepackage{xeCJK}
\IfFontExistsTF{Segoe UI}{\setmainfont{Segoe UI}}{\setmainfont{Arial}}
\IfFontExistsTF{Microsoft JhengHei}{\setCJKmainfont{Microsoft JhengHei}}{\IfFontExistsTF{Noto Sans CJK TC}{\setCJKmainfont{Noto Sans CJK TC}}{\setCJKmainfont{Noto Sans TC}}}
\IfFontExistsTF{Microsoft JhengHei}{\setCJKsansfont{Microsoft JhengHei}}{\IfFontExistsTF{Noto Sans CJK TC}{\setCJKsansfont{Noto Sans CJK TC}}{\setCJKsansfont{Noto Sans TC}}}
\xeCJKsetup{CJKglue={\hskip 0.045em plus 0.015em minus 0.005em}, CJKecglue={\hskip 0.08em plus 0.02em}}
\IfFileExists{../../../FontAwesome.otf}{\newfontfamily\FA[Path=../../../]{FontAwesome.otf}}{\newfontfamily\FA{Segoe UI Symbol}}
\newcommand{\faIcon}[1]{{\FA\symbol{#1}}}
\usetikzlibrary{arrows.meta,positioning,calc}
\definecolor{paper}{HTML}{FFFDF8}
\definecolor{navy}{HTML}{0F172A}
\definecolor{muted}{HTML}{334155}
\definecolor{line}{HTML}{CBD5E1}
\definecolor{green}{HTML}{00856A}
\definecolor{greensoft}{HTML}{DCFCE7}
\definecolor{gold}{HTML}{D97706}
\definecolor{goldsoft}{HTML}{FEF3C7}
\definecolor{red}{HTML}{DC2626}
\definecolor{redsoft}{HTML}{FEE2E2}
\begin{document}
\begin{tikzpicture}[x=1cm,y=1cm,>=Latex,line cap=round,line join=round]
\path[fill=paper] (0,0) rectangle (16,9.6);
\node[font=\bfseries\fontsize{18}{23}\selectfont,text=navy,align=center,text width=14.5cm] at (8,8.82) {沒有行動日誌，就沒有責任};
\node[font=\fontsize{10.3}{13.8}\selectfont,text=muted,align=center,text width=13.2cm] at (8,7.94) {Agent 不是只交結果；它要交出自己走過哪幾步、碰過哪些資料、在哪裡停下。};

\draw[rounded corners=12pt,fill=white,draw=line,line width=1pt] (1.05,2.05) rectangle (14.95,7.25);
\draw[rounded corners=12pt,fill=navy,draw=navy,line width=1pt] (1.05,6.45) rectangle (14.95,7.25);
\node[font=\bfseries\fontsize{10}{12}\selectfont,text=white,align=center,text width=2.2cm] at (2.35,6.85) {時間};
\node[font=\bfseries\fontsize{10}{12}\selectfont,text=white,align=center,text width=2.6cm] at (5.0,6.85) {動作};
\node[font=\bfseries\fontsize{10}{12}\selectfont,text=white,align=center,text width=2.8cm] at (8.0,6.85) {資料來源};
\node[font=\bfseries\fontsize{10}{12}\selectfont,text=white,align=center,text width=2.4cm] at (10.9,6.85) {結果};
\node[font=\bfseries\fontsize{10}{12}\selectfont,text=white,align=center,text width=2.3cm] at (13.3,6.85) {狀態};

\foreach \y in {5.55,4.65,3.75,2.85}{\draw[draw=line,line width=.7pt] (1.35,\y) -- (14.65,\y);}
\node[font=\fontsize{8.8}{11}\selectfont,text=muted,align=center,text width=2.1cm] at (2.35,6.0) {09:12};
\node[font=\fontsize{8.8}{11}\selectfont,text=navy,align=center,text width=2.55cm] at (5.0,6.0) {讀取會議附件};
\node[font=\fontsize{8.8}{11}\selectfont,text=muted,align=center,text width=2.75cm] at (8.0,6.0) {雲端資料夾};
\node[font=\fontsize{8.8}{11}\selectfont,text=navy,align=center,text width=2.35cm] at (10.9,6.0) {找到三份檔案};
\node[font=\bfseries\fontsize{8.6}{10}\selectfont,text=green,align=center,text width=2.2cm] at (13.3,6.0) {可繼續};

\node[font=\fontsize{8.8}{11}\selectfont,text=muted,align=center,text width=2.1cm] at (2.35,5.1) {09:18};
\node[font=\fontsize{8.8}{11}\selectfont,text=navy,align=center,text width=2.55cm] at (5.0,5.1) {補缺漏資料};
\node[font=\fontsize{8.8}{11}\selectfont,text=muted,align=center,text width=2.75cm] at (8.0,5.1) {外部網頁};
\node[font=\fontsize{8.8}{11}\selectfont,text=navy,align=center,text width=2.35cm] at (10.9,5.1) {來源不一致};
\node[font=\bfseries\fontsize{8.6}{10}\selectfont,text=gold,align=center,text width=2.2cm] at (13.3,5.1) {需標示};

\node[font=\fontsize{8.8}{11}\selectfont,text=muted,align=center,text width=2.1cm] at (2.35,4.2) {09:24};
\node[font=\fontsize{8.8}{11}\selectfont,text=navy,align=center,text width=2.55cm] at (5.0,4.2) {修改摘要草稿};
\node[font=\fontsize{8.8}{11}\selectfont,text=muted,align=center,text width=2.75cm] at (8.0,4.2) {副本};
\node[font=\fontsize{8.8}{11}\selectfont,text=navy,align=center,text width=2.35cm] at (10.9,4.2) {產生差異表};
\node[font=\bfseries\fontsize{8.6}{10}\selectfont,text=green,align=center,text width=2.2cm] at (13.3,4.2) {可審查};

\node[font=\fontsize{8.8}{11}\selectfont,text=muted,align=center,text width=2.1cm] at (2.35,3.3) {09:31};
\node[font=\fontsize{8.8}{11}\selectfont,text=navy,align=center,text width=2.55cm] at (5.0,3.3) {寄給外部收件人};
\node[font=\fontsize{8.8}{11}\selectfont,text=muted,align=center,text width=2.75cm] at (8.0,3.3) {正式信箱};
\node[font=\fontsize{8.8}{11}\selectfont,text=navy,align=center,text width=2.35cm] at (10.9,3.3) {不可逆};
\node[font=\bfseries\fontsize{8.6}{10}\selectfont,text=red,align=center,text width=2.2cm] at (13.3,3.3) {停下等人};

\draw[rounded corners=9pt,fill=redsoft,draw=red,line width=1pt] (5.7,.75) rectangle (10.3,1.55);
\node[font=\bfseries\fontsize{9.8}{12}\selectfont,text=navy,align=center,text width=4.1cm] at (8,1.15) {日誌的目的不是留痕，是讓錯誤能被追問。};
\end{tikzpicture}
\end{document}
